\chapter{Introducción}
\label{cap:introduccion}

\section{El problema y su contexto}

En los sistemas pensionales de capitalización individual, la renta vitalicia
inmediata permite transferir a una compañía de seguros el capital destinado a
financiar la pensión, a cambio de los pagos que correspondan mientras el afiliado
o sus beneficiarios conserven el derecho\footnote{El capital disponible puede incluir el saldo de la cuenta de ahorro individual y el bono pensional, cuando sea aplicable. Su comparación con una reserva requiere precisar fecha y base de valoración. La prima única del contrato puede incorporar componentes de tarificación que no coinciden con la reserva modelada.}.
En Colombia, la modalidad está regulada por la Ley 100 de
1993\footnote{Congreso de la República de Colombia, \emph{Ley 100 de 1993}, artículos 64 (requisitos para obtener la pensión de vejez), 74 (beneficiarios de la pensión de sobrevivientes), 79 (modalidades de las pensiones), 80 (renta vitalicia inmediata) y 81 (retiro programado). Los artículos 35 y 14 regulan la pensión mínima y el reajuste de las mesadas, respectivamente.} y por la reglamentación de la Superintendencia
Financiera.

Desde el punto de vista de la aseguradora, el contrato es una obligación de largo
plazo cuyo valor depende de variables demográficas (supervivencia del grupo
familiar), financieras (descuento e inflación) y regulatorias (indexación, piso
pensional y auxilio funerario). En el modelo de este trabajo, la equivalencia entre
un capital objetivo y el valor presente esperado de las obligaciones incluidas se
escribe como
\begin{equation}
\label{eq:problema-central}
  f(M) \;=\; \VP(M) - \KA \;=\; 0,\qquad M\ge\SMLV,
\end{equation}
donde $\KA>0$ es el capital objetivo, $M$ la mesada inicial y $\VP(M)$ la reserva
calculada por el cotizador. La igualdad exige bases de valoración homogéneas y no
identifica por sí sola esa reserva con una prima comercial completa.

La función combina supervivencias tabuladas, indexación, un piso expresado como
máximo y un auxilio funerario con recortes, sobre un horizonte de
\datoNMeses{} periodos mensuales. Aunque en cada tramo afín la inversión es
algebraica, localizar el tramo pertinente y controlar los casos de frontera
justifican estudiar los métodos numéricos que resuelven la ecuación.

Este trabajo pregunta qué métodos resultan adecuados para este cotizador y bajo
qué condiciones. La respuesta exige caracterizar primero la estructura de
$\VP$, y distinguir las propiedades demostradas del desempeño observado en la
muestra. El alcance se limita al modelo implementado, con las convenciones y
limitaciones actuariales explicitadas en los \cref{cap:modelo,cap:variables}.

\section{Motivación práctica}

La cotización repetida y la proyección anual del retiro programado requieren
resolver múltiples problemas de raíz. Cada llamada al motor construye flujos y
probabilidades y los descuenta sobre el horizonte de valoración; ese trabajo es
distinto de las operaciones escalares del método iterativo. Comparar solo el
número de iteraciones puede ocultar diferencias importantes: una derivada por
diferencias finitas requiere valoraciones adicionales, mientras que el cálculo
analítico conjunto puede reutilizar los vectores de una misma llamada.

La comparación del \cref{cap:resultados} contabiliza las llamadas efectivas al
motor, incluida la preparación del escenario, y las complementa con tiempos de
ejecución. Esa contabilidad permite evaluar el beneficio de la derivada exacta sin
suponer que su costo marginal sea nulo ni atribuir todo el costo de una cotización
al bucle del solucionador.

Hay además una razón de exactitud. El residuo se mide en unidades de capital,
pero el resultado que interesa se expresa en unidades de mesada. Una cota inferior
positiva de las pendientes relaciona ambos errores; cerca del piso pensional esa
relación puede ser menos favorable. El \cref{cap:metodos} establece criterios de
parada que distinguen estas unidades y el \cref{cap:retiro} estudia su propagación
entre años.

\section{Pertinencia del problema}
\label{sec:pertinencia}

La contribución académica consiste en analizar una función de valoración concreta
y relacionar su estructura con la elección y evaluación de métodos conocidos. La
resolución de una ecuación escalar adquiere aquí interés por la interacción de
las garantías con la inversión de la reserva, por los casos de capital
insuficiente y por el acoplamiento temporal del retiro programado.

\paragraph{Estructura de la función.}
Los modelos clásicos de anualidades sin garantías conducen a una valoración
lineal en el beneficio \citep{bowers1997,dickson2020,promislow2015}. Los máximos
y recortes del cotizador estudiado modifican esa estructura. En el
\cref{cap:estructura} se demuestra que:
\begin{enumerate}[label=(\roman*),leftmargin=*,nosep]
  \item la extensión de $\VP$ a $M\ge0$ es constante en $[0,\SMLV]$
        (\cref{thm:region-plana}); en el dominio admisible $M\ge\SMLV$ no hay
        solución si $\KA<\VP(\SMLV)$ y hay una única solución si
        $\KA\ge\VP(\SMLV)$, incluida $M=\SMLV$ en el caso de igualdad;
  \item el auxilio funerario puede destruir la convexidad bajo las condiciones
        del \cref{thm:no-convexidad}, por lo que no procede invocar sin
        comprobación la convergencia monótona de Newton para funciones convexas;
  \item la función es afín a trozos, con umbrales candidatos caracterizados
        explícitamente y agrupados en tres racimos geométricos para los
        parámetros considerados (\cref{thm:quiebres}).
\end{enumerate}
En la extensión $M\ge0$, el caso $\KA=\VP(\SMLV)$ sí tiene un intervalo de
raíces; esa pérdida de inyectividad no debe confundirse con la unicidad en el
dominio admisible.

\paragraph{Análisis no suave y consecuencias algorítmicas.}
Las derivadas laterales y los subdiferenciales de Clarke y de Bouligand
\citep{clarke1983,qi1993} describen los puntos de quiebre. La semisuavidad fuerte
y la $BD$\nobreakdash-regularidad en una solución interior permiten aplicar la
teoría local de Newton semisuave \citep{qi1993,facchinei2003}
(\crefrange{thm:semisuavidad}{thm:bd-regularidad}). La inversión de un tramo con
su pendiente exacta da la raíz en un paso si su clausura contiene también la
solución. La garantía global requiere condiciones adicionales, como las de la
salvaguarda de encaje con bisección periódica; no se deduce de la sola linealidad
a trozos.

\paragraph{Resultados aplicados verificables.}
El trabajo presenta la expresión del subdiferencial para este modelo, las
condiciones de terminación, el sesgo de las diferencias finitas en los quiebres y
una cota condicional para el estado de retiro programado
(\cref{sec:aportes}). La novedad se sitúa en su desarrollo e integración para el
cotizador examinado, y no en la introducción de una teoría general nueva de
funciones afines a trozos. Los experimentos contrastan implicaciones de esos
resultados; no sustituyen sus demostraciones ni validan las bases actuariales.

\paragraph{Delimitación.}
Se estudia una ecuación escalar y una recursión anual con un conjunto explícito de
convenciones. El instrumental comprende análisis no suave, convergencia de
métodos iterativos y condiciones para intercambiar esperanza y suma de flujos
(\cref{thm:tonelli}). Los \datoNMetodos{} métodos comparados son conocidos. El
aporte central es explicar cómo aprovechar la estructura de este problema y cómo
medir el trabajo necesario para resolverlo con precisión controlada.

\begin{criticism}[Sobre el alcance]
Una reserva bien resuelta numéricamente puede seguir apoyándose en tablas o
parámetros insuficientemente documentados. La derivada exacta y el encaje permiten
controlar el error del solucionador, mientras que la trazabilidad demográfica,
la calibración financiera y la conciliación con cotizaciones requieren evidencia
propia. Esta separación es necesaria para interpretar tanto las mesadas de renta
vitalicia como las trayectorias ilustrativas de retiro programado.
\end{criticism}

\section{Estado del arte}
\label{sec:estado-arte}

La revisión conecta tres frentes.

\paragraph{Valoración actuarial de rentas contingentes.}
\citet{bowers1997,dickson2020,gerber1997,promislow2015} desarrollan el valor
presente actuarial y la equivalencia en anualidades sobre una y varias vidas. En
los modelos clásicos con beneficios proporcionales a $M$, la inversión es lineal.
El interés de esta aplicación está en conservar explícitos los máximos y recortes
del cotizador para estudiar su inversión y sus casos de frontera.

\paragraph{Garantías y análisis no suave.}
Las garantías de productos de seguro introducen opcionalidad
\citep[cap.~11]{dickson2020}. La teoría convexa y no suave de
\citet{rockafellar1970,boyd2004,clarke1983} proporciona herramientas para
analizar las operaciones de máximo, mínimo y suma que intervienen en este modelo.
La presencia de garantías no determina por sí sola la forma de la función: esta
depende también de la indexación, los pesos y las demás hipótesis de valoración.

\paragraph{Métodos iterativos y su análisis.}
Los tratamientos de \citet{burden2016,atkinson1989,stoer2002,quarteroni2007}
fundamentan los métodos escalares. El catálogo incluye alternativas con encaje de
\citet{brent1973,ridders1979,dowell1972} e ITP \citep{oliveira2021}. Para el
caso no suave se emplean \citet{qi1993,facchinei2003}. El índice de eficiencia de
\citet{ostrowski1966} se interpreta junto con el costo efectivo, y no como un
sustituto de su medición.

\paragraph{La contribución en esta intersección.}
El trabajo combina estas herramientas para caracterizar el cotizador estudiado,
construir sus derivadas laterales e instrumentar una comparación de métodos.
La revisión presentada delimita ese aporte aplicado; no acredita ausencia de
trabajos similares en toda la literatura ni prioridad universal.

\section{Pregunta de investigación, objetivos e hipótesis}

\subsection*{Pregunta de investigación}
\begin{quote}
\textbf{¿Qué propiedades estructurales de la función de valoración del cotizador
estudiado determinan la elección de métodos iterativos, y cómo se comparan estos
en precisión, robustez y costo sobre los escenarios examinados?}
\end{quote}

La caracterización analítica sustenta la comparación: un resultado experimental
favorable debe interpretarse junto con las condiciones que permiten o impiden
generalizarlo.

\subsection*{Objetivo general}
Caracterizar la estructura analítica de la función de valoración del modelo
estudiado y evaluar, sobre esa base, métodos iterativos para resolver la ecuación
de cotización con precisión, robustez y costo documentados.

\subsection*{Objetivos específicos}
\begin{enumerate}[leftmargin=*]
  \item Formular la cotización como problema de raíz y establecer condiciones de
        existencia y unicidad en su dominio admisible. (\cref{cap:estructura})
  \item Caracterizar los umbrales candidatos y las propiedades de monotonía,
        convexidad, semisuavidad y condicionamiento. (\cref{cap:estructura})
  \item Obtener una expresión del subdiferencial e interpretarla actuarialmente.
        (\cref{cap:estructura})
  \item Implementar y analizar \datoNMetodos{} métodos de encaje, abiertos y no
        suaves, distinguiendo garantías locales y globales. (\cref{cap:metodos})
  \item Revisar el fundamento y la sensibilidad de las variables demográficas,
        financieras y regulatorias. (\cref{cap:variables})
  \item Extender la formulación al retiro programado y analizar la propagación del
        error en su estado acoplado. (\cref{cap:retiro})
  \item Comparar los métodos sobre escenarios sintéticos y muestras de pólizas,
        documentando preparación, evaluaciones y precisión. (\cref{cap:resultados})
\end{enumerate}

\subsection*{Hipótesis de trabajo}
\begin{description}[leftmargin=1.2cm,style=nextline]
  \item[H1] La valoración es afín a trozos en la mesada inicial, con un conjunto
        finito y explícitamente caracterizable de umbrales candidatos.
  \item[H2] La presencia de intervalos afines permite terminar en un paso cuando
        el punto de evaluación y la raíz pertenecen a la clausura del tramo cuya
        pendiente se selecciona.
  \item[H3] El número de iteraciones es insuficiente para comparar el costo; es
        necesario contar las llamadas efectivas al motor y medir el tiempo.
  \item[H4] La derivada analítica conjunta permite evitar las valoraciones
        auxiliares de las diferencias finitas; su beneficio efectivo depende de
        la instrumentación y de los escenarios de prueba.
\end{description}

H1 y H2 se estudian analíticamente; H3 y H4 se contrastan mediante la contabilidad
y las mediciones del \cref{cap:resultados}.

\section{Aportes del trabajo}
\label{sec:aportes}

\begin{aporte}
\textbf{Teóricos y de modelación.}
\begin{enumerate}[leftmargin=*,nosep]
  \item Región plana de la extensión y condición de capital mínimo en el dominio
        admisible (\cref{thm:region-plana,cor:capital-minimo}).
  \item Umbrales candidatos en tres racimos y condiciones de pérdida de convexidad
        (\cref{thm:quiebres,thm:no-convexidad}).
  \item Expresión del $B$-subdiferencial, interpretación marginal, semisuavidad y
        $BD$-regularidad interior
        (\cref{thm:subdiferencial,thm:semisuavidad,thm:bd-regularidad}).
  \item Condiciones de terminación de Newton y garantía de encaje para la variante
        salvaguardada (\cref{thm:terminacion-finita,thm:salvaguardado}).
  \item Sesgo de las diferencias finitas en quiebres y cota condicional del error
        en el estado de retiro programado
        (\cref{prop:sesgo-diferencias,thm:propagacion}).
\end{enumerate}

\textbf{Metodológicos y computacionales.}
\begin{enumerate}[leftmargin=*,nosep,resume]
  \item Implementación de \datoNMetodos{} métodos con interfaz común y
        contabilidad de llamadas completas al motor, memoria por ejecución y
        descripción de las trayectorias de error sin atribuir orden asintótico
        a una regresión de pocas iteraciones.
  \item Evaluación conjunta de reserva y derivadas laterales en el motor.
  \item Protocolo con raíces de referencia y escenarios que incluyen quiebres,
        regiones afines y capital insuficiente.
  \item Hallazgos y recomendaciones sobre la implementación
        (\cref{sec:hallazgos-codigo}).
  \item Aplicación interactiva para explorar el modelo y sus resultados.
\end{enumerate}
\end{aporte}

\section{Estructura del documento}

El \cref{cap:preliminares} reúne los preliminares de probabilidad, análisis no
suave y convergencia. El \cref{cap:modelo} formula el modelo y delimita las
convenciones del cotizador. El \cref{cap:estructura} desarrolla el análisis
estructural y el \cref{cap:metodos} los métodos y sus garantías. El
\cref{cap:variables} revisa las variables; el \cref{cap:retiro} estudia la
recursión de retiro programado. El \cref{cap:resultados} presenta los experimentos
y el \cref{cap:conclusiones} sus implicaciones y límites. Los apéndices recogen la
notación y el software. El proyecto identifica como repositorio del trabajo
\repositorio; la reproducción requiere la versión del código y los datos
asociados a cada ejecución.
